\documentclass[a4paper,10pt]{article}
\usepackage[brazilian]{babel}
\usepackage{lmodern}
\usepackage[T1]{fontenc}
\usepackage[left=0.5in,right=0.5in,top=0.45in,bottom=0.45in]{geometry}
\usepackage{titlesec}
\usepackage{enumitem}
\usepackage[hidelinks]{hyperref}
\usepackage{tabularx}

\pagestyle{empty}
\setlength{\parindent}{0pt}
\raggedbottom

\titleformat{\section}{\Large\bfseries}{}{0em}{}[\vspace{-2pt}\titlerule]
\titlespacing*{\section}{0pt}{12pt}{8pt}

\newcommand{\entry}[4]{%
  \begin{tabularx}{\textwidth}{@{}X r@{}}

    \textbf{\large #1} & \textbf{\large #2} \\
    \textit{#3} & \textit{#4} \\
  \end{tabularx}\vspace{1pt}}
\newenvironment{items}{\begin{itemize}[leftmargin=26pt,itemsep=3pt,parsep=0pt,topsep=0pt,partopsep=0pt,label=\textbullet]}{\end{itemize}\vspace{8pt}}

\begin{document}

\begin{center}
  {\LARGE\bfseries Sara Cerutti}\\[3pt]
  \small E-mail: \href{mailto:saracmuller@gmail.com}{saracmuller@gmail.com} \textbullet\ \href{https://linkedin.com/in/saramllr/}{linkedin.com/in/saramllr/}
\end{center}

\section{Product Designer | UX/UI | Pesquisa e Estratégia}
Product Designer com mais de 6 anos de experiência no design de produtos digitais complexos em e-commerce, educação e saúde, com atuação em pesquisa com usuários, product discovery, design de interação e UI. Trabalho de ponta a ponta, da pesquisa e da definição do problema ao design em alta fidelidade, à prototipação e à implementação. Tenho experiência em simplificar sistemas com múltiplos perfis, estruturar fundamentos de design e transformar pesquisa e dados de comportamento em decisões de produto. Trabalho perto dos times de Produto e Engenharia em ambientes de ritmo acelerado e uso IA em todo o meu fluxo para acelerar a síntese de pesquisa, a exploração, a prototipação, a documentação e a entrega.

\section{Experiência}

\entry{+A Educação}{Remoto}{Product Designer}{junho de 2021 - atual}
\begin{items}
  \item Liderei a pesquisa de discovery da reconstrução do GPA, um produto complexo com múltiplos perfis, combinando avaliação heurística, entrevistas com usuários, um questionário com 55 respostas e um piloto em produção. Os achados influenciaram diretamente as prioridades do roadmap, incluindo a decisão de priorizar o mobile no lugar de uma funcionalidade pedida pelo negócio.
  \item Redesenhei fluxos centrais das experiências de estudante, professor e coordenador como parte de uma reconstrução de produto de vários anos, resolvendo atritos e problemas estruturais do produto anterior. Conduzi pesquisas com usuários (entrevistas, testes de usabilidade e questionários), gerando insights para a melhoria contínua do produto.
  \item Desenhei o Artemya, uma plataforma de treino clínico para estudantes de Psicologia criada do zero, dos primeiros fluxos e modelos de interação à prototipação em alta fidelidade. Os testes de usabilidade orientaram iterações pontuais, preservando o atrito intencional nas interações essenciais para o aprendizado.
  \item Usei Microsoft Clarity e Google Tag Manager para investigar o comportamento dos usuários, identificar problemas de usabilidade e embasar decisões de design com evidências de comportamento.
  \item Contribuí com o design system e os fundamentos de design do time, documentando componentes, estados e padrões com foco em acessibilidade, consistência e reuso em escala.
  \item Criei fluxos de trabalho com IA para o processo de design, incluindo validação de requisitos e especificações automatizadas para o handoff com o desenvolvimento, transformando experimentos individuais em práticas repetíveis do time.
  \item Trabalhei lado a lado com Product Owners e Engenharia em discovery, design, implementação e iteração, equilibrando as necessidades dos usuários, as restrições técnicas e as prioridades do produto.
\end{items}

\entry{Iú Cosmetics}{Campo Bom, Brasil}{Digital Designer (E-commerce e Experiência)}{outubro de 2019 - março de 2021}
\begin{items}
  \item Melhorei a jornada de compra do e-commerce em navegação, páginas de produto e fluxos de campanha de uma marca de beleza com venda direta ao consumidor (DTC).
  \item Fui responsável pelas entregas de design e comunicação durante a migração completa da plataforma de e-commerce, coordenando produto, marketing e operações.
  \item Desenhei peças voltadas à conversão para listagens de produto, e-mail marketing e campanhas pagas, mantendo a consistência visual em todos os pontos de contato com o cliente.
  \item Produzi fotografia de produto e conteúdo visual estruturado para deixar os produtos mais claros no momento da decisão de compra.
  \item Colaborei com várias áreas em um ambiente DTC de ritmo acelerado, equilibrando necessidades do negócio, qualidade criativa e prazos de entrega.
\end{items}

\entry{Luedtke}{São Leopoldo, Brasil}{Designer Instrucional}{março de 2017 - dezembro de 2017}
\begin{items}
  \item Desenhei experiências de aprendizagem digital, incluindo cursos e-learning, conteúdo interativo, materiais visuais e vídeos.
  \item Criei cursos e-learning no Articulate Storyline e produzi ilustrações, animações e vídeos explicativos de apoio.
\end{items}

\entry{Agexcom}{São Leopoldo, Brasil}{Estagiária de Direção de Arte}{agosto de 2014 - junho de 2016}
\begin{items}
  \item Contribuí com campanhas digitais e experiências visuais, do conceito à entrega final.
  \item Desenvolvi ilustrações, motion graphics e conteúdo digital em diversos formatos.
\end{items}

\section{Habilidades}
\begin{items}
  \item \textbf{Técnicas:} Experiência do Usuário (UX), Interface do Usuário (UI), UX Research, Pesquisa com Usuários, Entrevistas com Usuários, Testes de Usabilidade, Questionários, Análise de Dados, Product Discovery, Design de Interação, Arquitetura da Informação, Fluxos de Usuário, Jornada do Usuário, Wireframes, Prototipação, Design Systems, Padrões de UI, Heurísticas de Usabilidade, Acessibilidade, Design Inclusivo, UX Writing, Content Design, Storytelling, Design Estratégico, Metodologias Ágeis (Scrum, Kanban), Gestão de Stakeholders, Decisões Orientadas por Dados, Testes A/B, Analytics, Melhoria Contínua, EdTech, E-learning, LXP, E-commerce
  \item \textbf{Ferramentas:} Figma, Adobe Creative Cloud (Photoshop, Illustrator, After Effects), Google Analytics, Hotjar, Microsoft Clarity, ChatGPT, Claude
  \item \textbf{Competências:} Comunicação clara e estruturada, Pensamento crítico, Colaboração entre áreas, Gestão de stakeholders, Pensamento sistêmico, Empatia aplicada à experiência do usuário, Tradução de problemas complexos em soluções simples, Design centrado no comportamento
  \item \textbf{Idiomas:} Português (nativo), Inglês (fluente), Espanhol (básico)
\end{items}

\section{Formação}

\entry{Universidade do Vale do Rio dos Sinos (UNISINOS)}{Remoto}{MBA em Design Estratégico (em andamento)}{2026 - 2027}
\vspace{6pt}

\entry{Pontifícia Universidade Católica do Paraná (PUC-PR)}{Remoto}{Pós-graduação em Gamificação, Engajamento e Marketing Digital}{setembro de 2021 - setembro de 2022}
\vspace{6pt}

\entry{Universidade do Vale do Rio dos Sinos (UNISINOS)}{São Leopoldo, RS}{Bacharelado em Comunicação Digital}{janeiro de 2014 - dezembro de 2017}

\end{document}
